% articulo-12-notas.tex - Notas al pie.
% (footmisc + \@makefntext)
%
% FILETE SEPARADOR (Chicago § 16.63)
%
%   Comportamiento por defecto ([norule,splitrule]):
%     - Página con nota normal o inicio de nota larga: sin filete.
%     - Página de continuación (nota de la página anterior): filete sin sangría.
%   Esto es lo que Chicago indica como obligatorio.
%
%   Variables YAML:
%     filete-pie:   true  →  filete sin sangría en todas las páginas con notas.
%     filete-nunca: true  →  sin filete en ningún caso (suprime también el de continuación).
%     (sin declarar)      →  filete solo en páginas de continuación (por defecto).
%
%   Mecanismo: [norule] fija \pagefootnoterule a vacío; [splitrule] usa
%   \insertpenalties como flag para distinguir página normal de continuación.
%   Con filete-pie se redefine \pagefootnoterule con las mismas dimensiones
%   que \splitfootnoterule (0.6 pt × 1.5 in, sin sangría).
%
% FORMATO DEL PIE (Chicago § 16.63)
%   - Número (superíndice) sangrado \ParIndent desde el margen.
%   - Medio cuadratín (0.5em) entre número y texto.
%   - Párrafos adicionales: sangría \ParIndent en primera línea.
%   - Listas: mismo formato que en el cuerpo (rayas, \ParIndent, etc.)
%
%   Variables YAML:
%     (sin declarar)          →  primera línea de cada nota sangrada (por defecto)
%     notas-sin-sangria: true →  primera línea de cada nota sin sangría
%
% FAMILIA TIPOGRÁFICA DEL NUMERADOR (marca en cuerpo y en pie)
%   El texto de la nota va siempre en la familia principal del documento (serif
%   si el cuerpo es serif), siguiendo la convención tipográfica habitual.
%   El numerador, como elemento del aparato de referencia, puede ir en sans
%   cuando titulos-sans: true, en coherencia con encabezados, folios e índice.
%
%   Variables YAML:
%     titulos-sans: true  →  numerador (superíndice en cuerpo y en pie) en sans.
%     (sin declarar)      →  numerador en la familia principal (por defecto).
%
% La redefinición de \@makefntext va después de que footmisc ejecute
% \ProcessOptions (que es donde footmisc redefine \@makefntext).
%
% \@parboxrestore (llamado por footmisc justo antes de \@makefntext) resetea
% \parindent a 0pt. Por eso se restaura \parindent=\ParIndent DENTRO de
% \@makefntext, antes de procesar #1, para que los párrafos adicionales y
% las listas del cuerpo de la nota funcionen correctamente.
%
% Se resetean también \@enumdepth e \@itemdepth a 0. Cuando el footnote
% aparece dentro de un \item, enumitem lleva ya un nivel de anidamiento
% acumulado; sin el reset, la primera lista dentro de la nota se trataría
% como nivel 2 (boliche) en lugar de nivel 1 (raya).

$if(filete-nunca)$
\usepackage[norule]{footmisc}
$else$$if(filete-pie)$
\usepackage[splitrule]{footmisc}
$else$
\usepackage[norule,splitrule]{footmisc}
$endif$$endif$

\renewcommand{\footnotesize}{\small}

% Filetes separadores: 0.6 pt de grosor, 1.5 pulgadas de longitud, sin sangría.
% \splitfootnoterule: páginas de continuación (siempre activo salvo filete-nunca).
% \pagefootnoterule:  páginas de nota normal (solo con filete-pie).
\makeatletter
\newcommand{\@footrule}{%
  \kern-3\p@
  \hrule \@width 1.5in \@height 0.6pt
  \kern 2.6\p@
}
\let\splitfootnoterule\@footrule
$if(filete-pie)$
\let\pagefootnoterule\@footrule
$endif$
\makeatother

% Marca en el cuerpo del texto: superíndice en negro.
% hyperref colorea por defecto con linkcolor (azul); se sobreescribe aquí.
% Con titulos-sans: true, el numerador va en sans como el resto del aparato
% navegacional; sin declarar, hereda la familia principal.
\renewcommand{\thefootnote}{%
  \textcolor{black}{$if(titulos-sans)${\sffamily\arabic{footnote}}$else$\arabic{footnote}$endif$}%
}

% Formato del pie
% (sin declarar)             → número sangrado + medio cuadratín + texto.
% notas-sin-sangria: true    → la primera línea no lleva sangría.
% nota-numeracion-linea true → «1. Texto» en línea, sin superíndice
\makeatletter
\renewcommand{\@makefntext}[1]{%
  \parindent=\ParIndent%             % restaurar (footmisc + @parboxrestore lo ponen a 0pt)
  \@enumdepth=0\relax%               % resetear profundidad de enumerate
  \@itemdepth=0\relax%               % resetear profundidad de itemize
  \noindent
$if(notas-sin-sangria)$%
$else$%
  \hspace*{\ParIndent}%
$endif$%
$if(nota-numeracion-linea)$
  $if(titulos-sans)${\sffamily\arabic{footnote}}$else$\arabic{footnote}$endif$.
$else$%
  \@makefnmark%
$endif$%
  \hspace{0.5em}#1%
}
\makeatother
